%!TEX root = iopjournal.tex
%
% Injecting real concept drift into the TEP
%
% Erster eigener Beitrag. Die IDV(13)-Analyse steht bewusst hier und nicht im
% Background: sie ist ein eigener Befund und begruendet, dass es IDV(29) braucht.
%
% Inhalt (Vorschlag):
%   - IDV(13): r1f/r2f via tesub5/tesub8, stationaer und ohne Trend -> kein CD
%   - Wahl des Angriffspunkts: Anlagenphysik wird ausgeregelt und bleibt im
%     Messvektor sichtbar -> nur Covariate Shift; die Sensorik liegt ausserhalb
%     dieses Wirkungskreises
%   - IDV(29): multiplikativer Kalibrierfehler 1 + kappa_m * c(t) auf XMEAS(18);
%     Begruendung ueber die drei Eigenschaften (aufgezeichnet, kein Regelkreis,
%     nicht in der Kostenformel) -> Verschiebung von P(y|X)
%   - Driftsignal c(t): zustandslos als geschlossener Ausdruck in t, weil ode45
%     mit veraenderlicher Schrittweite rechnet; Segmentraster, Sudden und
%     Incremental Drift in einem Signal
%   - Eingriff im Quelltext (Listing aus temexd_mod.c)
%
% Aus Abschnitt 3 offen:
%   - Einordnung explizit machen: real concept drift, aber kein actual drift --
%     P(X) wandert mit, weil XMEAS(18) Teil des Merkmalsvektors ist
%   - c(t) auf die Musterachse aus 3.1 abbilden: inkrementell im Segment,
%     sudden am Ruecksprung, recurring durch die Rueckkehr auf den Nominalwert;
%     gradual kommt im Strom nicht vor
%
% Vgl. Dissertation: sec:tep_drift
%

\section{Injecting real concept drift into the Tennessee Eastman process} \label{sec:drift_injection}
%
Among the disturbances of the catalog, \texttt{IDV(13)} appears to come closest to the purpose of this work, since it is labeled a \emph{slow drift} of the reaction kinetics~\cite{downs.1993}. Once it is activated, the simulation multiplies the rates of the two primary reactions, which form the liquid products, by the time-varying factors \texttt{r1f} and \texttt{r2f}. At random intervals, each factor draws a new target value uniformly from a bounded interval of $[0.75, 1.25]$ around its nominal value of $1$ and approaches it by cubic interpolation, so that its course is smooth but strictly bounded. The memory of this course therefore extends no further than the next target value, which lies only a few hours ahead. The resulting fluctuation causes a one-time expansion of the distribution of the affected process variables once the disturbance is activated. However, it does not shift the distribution any further, and the process remains stationary and without a trend. Since a concept drift according to~\eqref{eq:concept_drift} requires the joint distribution to change over time, \texttt{IDV(13)} does not constitute a concept drift despite its label, but a disturbance of stationary character. \par
%
The simulation is therefore extended by a further disturbance, \texttt{IDV(29)}, which is activated and scaled through the same interface as its other entries. Regarding the place where the drift may occur, an impact on the physics of the system—such as the reaction kinetics discussed in \texttt{IDV(13)}—seems plausible, since it represents an actual change in the process behavior. The TEP, however, is fully controlled, and its control scheme converts any persistent disturbance into a lasting change of the manipulated variables, which are recorded along with the measured ones. The operating costs, in turn, are computed by the simulator from this very measurement vector. A change in the system’s physics thus remains visible in the observed variables, and every quantity the simulator derives therefrom is calculated based on the changed values themselves. Thus, its relationship to these variables is preserved. What shifts is the position of the data within the domain of the observed variables, which, according to~\eqref{eq:covariate_shift}, constitutes a covariate shift. \par
%
However, the sensors that provide the measured values are located outside this chain of effects. In industrial practice, sensors are subject to fouling and aging, while maintenance and calibration intervals are limited. It is therefore assumed that a sensor provides increasingly distorted values over a period of weeks, even though the process itself remains unchanged. \texttt{IDV(29)} reproduces this kind of sensor deterioration on the stripper temperature \texttt{XMEAS(18)}, highlighted in red in figure~\ref{fig:tep_flowsheet}, which is chosen for three properties. It is recorded by the simulation, it is not part of any control loop, and it does not enter the calculation of the operating costs. The drift therefore distorts only the value that is output, while the process value, and thus the behavior of the plant, remains unaffected. \par
%
\texttt{IDV(29)} realizes this deterioration as a multiplicative calibration error and scales the output value of \texttt{XMEAS(18)} by the factor $1 + \kappa_m\, c(t)$, in which $\kappa_m$ denotes the drift coefficient and $c(t)$ the drift signal. Since the factor acts multiplicatively, the absolute error of the measurement depends on the operating point at which the plant is run. The operating costs, taken as the target $y$, continue to follow the unchanged process, whereas the recorded temperature departs increasingly from it. Over the course of the drift, the same observed value thus corresponds to different operating costs, so that the posterior $P(y \mid X)$ shifts and the drift can be classified according to~\eqref{eq:real_concept_drift} as a real concept drift. However, as the distorted measurement itself is one of the recorded variables, the input distribution shifts as well, so that the drift does not constitute an actual drift. \par
%
The drift signal $c(t)$ determines how this error evolves over time and is modeled according to a sensor found in real-world applications. Between maintenance intervals, fouling and aging cause the error to change continuously, while recalibration resets it immediately. The simulated period is therefore divided into segments whose nominal duration $T_{\mathrm{seg}}$ represents the calibration interval. Since a maintenance date varies in practice, the inner segment boundaries $\tau_1 < \tau_2 < \cdots$ are displaced randomly from this grid by a fraction $\gamma_\tau$ of the segment duration. Within the $i$-th segment, the signal falls linearly from a start level $c_i^{+}$, drawn from an upper band, to a target level $c_i^{-}$, drawn from a lower band,
%
\begin{equation} \label{eq:drift_signal}
    c(t) = c_i^{+} + \frac{c_i^{-} - c_i^{+}}{\tau_i - \tau_{i-1}}\,(t - \tau_{i-1}),
    \qquad t \in \bigl[\tau_{i-1},\, \tau_i\bigr),
\end{equation}
%
and at each boundary it jumps to the start level of the following segment. Figure~\ref{fig:drift_signal} shows the resulting course of $c(t)$ over the simulated period. Each segment appears as a separate line that falls from the upper to the lower shaded band, and the gaps between the lines mark the sudden drifts at the segment boundaries. This sawtooth combines most of the patterns introduced in section~\ref{subsec:concept_drift}. Within a segment, the concept passes through intermediate states, which constitutes an incremental drift, whereas at a boundary it is replaced at once, which constitutes a sudden drift. As the signal crosses zero in every segment, where the measurement is undistorted, the original concept is revisited repeatedly, so that the drift is recurring as well. A gradual drift, in which two concepts alternate, does not occur. \par
%
\begin{figure}[tb]
    \centering
    %% Creator: Matplotlib, PGF backend
%%
%% To include the figure in your LaTeX document, write
%%   \input{<filename>.pgf}
%%
%% Make sure the required packages are loaded in your preamble
%%   \usepackage{pgf}
%%
%% Also ensure that all the required font packages are loaded; for instance,
%% the lmodern package is sometimes necessary when using math font.
%%   \usepackage{lmodern}
%%
%% Figures using additional raster images can only be included by \input if
%% they are in the same directory as the main LaTeX file. For loading figures
%% from other directories you can use the `import` package
%%   \usepackage{import}
%%
%% and then include the figures with
%%   \import{<path to file>}{<filename>.pgf}
%%
%% Matplotlib used the following preamble
%%   \def\mathdefault#1{#1}
%%   \everymath=\expandafter{\the\everymath\displaystyle}
%%   
%%   \ifdefined\pdftexversion\else  % non-pdftex case.
%%     \usepackage{fontspec}
%%   \fi
%%   \makeatletter\@ifpackageloaded{underscore}{}{\usepackage[strings]{underscore}}\makeatother
%%
\begingroup%
\makeatletter%
\begin{pgfpicture}%
\pgfpathrectangle{\pgfpointorigin}{\pgfqpoint{6.066574in}{1.709658in}}%
\pgfusepath{use as bounding box, clip}%
\begin{pgfscope}%
\pgfsetbuttcap%
\pgfsetmiterjoin%
\definecolor{currentfill}{rgb}{1.000000,1.000000,1.000000}%
\pgfsetfillcolor{currentfill}%
\pgfsetlinewidth{0.000000pt}%
\definecolor{currentstroke}{rgb}{1.000000,1.000000,1.000000}%
\pgfsetstrokecolor{currentstroke}%
\pgfsetdash{}{0pt}%
\pgfpathmoveto{\pgfqpoint{0.000000in}{0.000000in}}%
\pgfpathlineto{\pgfqpoint{6.066574in}{0.000000in}}%
\pgfpathlineto{\pgfqpoint{6.066574in}{1.709658in}}%
\pgfpathlineto{\pgfqpoint{0.000000in}{1.709658in}}%
\pgfpathlineto{\pgfqpoint{0.000000in}{0.000000in}}%
\pgfpathclose%
\pgfusepath{fill}%
\end{pgfscope}%
\begin{pgfscope}%
\pgfsetbuttcap%
\pgfsetmiterjoin%
\definecolor{currentfill}{rgb}{1.000000,1.000000,1.000000}%
\pgfsetfillcolor{currentfill}%
\pgfsetlinewidth{0.000000pt}%
\definecolor{currentstroke}{rgb}{0.000000,0.000000,0.000000}%
\pgfsetstrokecolor{currentstroke}%
\pgfsetstrokeopacity{0.000000}%
\pgfsetdash{}{0pt}%
\pgfpathmoveto{\pgfqpoint{0.665489in}{0.450111in}}%
\pgfpathlineto{\pgfqpoint{5.333796in}{0.450111in}}%
\pgfpathlineto{\pgfqpoint{5.333796in}{1.609658in}}%
\pgfpathlineto{\pgfqpoint{0.665489in}{1.609658in}}%
\pgfpathlineto{\pgfqpoint{0.665489in}{0.450111in}}%
\pgfpathclose%
\pgfusepath{fill}%
\end{pgfscope}%
\begin{pgfscope}%
\pgfpathrectangle{\pgfqpoint{0.665489in}{0.450111in}}{\pgfqpoint{4.668307in}{1.159547in}}%
\pgfusepath{clip}%
\pgfsetbuttcap%
\pgfsetmiterjoin%
\definecolor{currentfill}{rgb}{0.301961,0.301961,0.301961}%
\pgfsetfillcolor{currentfill}%
\pgfsetfillopacity{0.120000}%
\pgfsetlinewidth{0.000000pt}%
\definecolor{currentstroke}{rgb}{0.301961,0.301961,0.301961}%
\pgfsetstrokecolor{currentstroke}%
\pgfsetstrokeopacity{0.120000}%
\pgfsetdash{}{0pt}%
\pgfpathmoveto{\pgfqpoint{0.665489in}{1.150671in}}%
\pgfpathlineto{\pgfqpoint{5.333796in}{1.150671in}}%
\pgfpathlineto{\pgfqpoint{5.333796in}{1.513029in}}%
\pgfpathlineto{\pgfqpoint{0.665489in}{1.513029in}}%
\pgfpathlineto{\pgfqpoint{0.665489in}{1.150671in}}%
\pgfpathclose%
\pgfusepath{fill}%
\end{pgfscope}%
\begin{pgfscope}%
\pgfpathrectangle{\pgfqpoint{0.665489in}{0.450111in}}{\pgfqpoint{4.668307in}{1.159547in}}%
\pgfusepath{clip}%
\pgfsetbuttcap%
\pgfsetmiterjoin%
\definecolor{currentfill}{rgb}{0.301961,0.301961,0.301961}%
\pgfsetfillcolor{currentfill}%
\pgfsetfillopacity{0.120000}%
\pgfsetlinewidth{0.000000pt}%
\definecolor{currentstroke}{rgb}{0.301961,0.301961,0.301961}%
\pgfsetstrokecolor{currentstroke}%
\pgfsetstrokeopacity{0.120000}%
\pgfsetdash{}{0pt}%
\pgfpathmoveto{\pgfqpoint{0.665489in}{0.546740in}}%
\pgfpathlineto{\pgfqpoint{5.333796in}{0.546740in}}%
\pgfpathlineto{\pgfqpoint{5.333796in}{0.909098in}}%
\pgfpathlineto{\pgfqpoint{0.665489in}{0.909098in}}%
\pgfpathlineto{\pgfqpoint{0.665489in}{0.546740in}}%
\pgfpathclose%
\pgfusepath{fill}%
\end{pgfscope}%
\begin{pgfscope}%
\pgfsetbuttcap%
\pgfsetroundjoin%
\definecolor{currentfill}{rgb}{0.000000,0.000000,0.000000}%
\pgfsetfillcolor{currentfill}%
\pgfsetlinewidth{0.803000pt}%
\definecolor{currentstroke}{rgb}{0.000000,0.000000,0.000000}%
\pgfsetstrokecolor{currentstroke}%
\pgfsetdash{}{0pt}%
\pgfsys@defobject{currentmarker}{\pgfqpoint{0.000000in}{-0.048611in}}{\pgfqpoint{0.000000in}{0.000000in}}{%
\pgfpathmoveto{\pgfqpoint{0.000000in}{0.000000in}}%
\pgfpathlineto{\pgfqpoint{0.000000in}{-0.048611in}}%
\pgfusepath{stroke,fill}%
}%
\begin{pgfscope}%
\pgfsys@transformshift{0.665489in}{0.450111in}%
\pgfsys@useobject{currentmarker}{}%
\end{pgfscope}%
\end{pgfscope}%
\begin{pgfscope}%
\definecolor{textcolor}{rgb}{0.000000,0.000000,0.000000}%
\pgfsetstrokecolor{textcolor}%
\pgfsetfillcolor{textcolor}%
\pgftext[x=0.665489in,y=0.352889in,,top]{\color{textcolor}{\rmfamily\fontsize{8.000000}{9.600000}\selectfont\catcode`\^=\active\def^{\ifmmode\sp\else\^{}\fi}\catcode`\%=\active\def%{\%}0}}%
\end{pgfscope}%
\begin{pgfscope}%
\pgfsetbuttcap%
\pgfsetroundjoin%
\definecolor{currentfill}{rgb}{0.000000,0.000000,0.000000}%
\pgfsetfillcolor{currentfill}%
\pgfsetlinewidth{0.803000pt}%
\definecolor{currentstroke}{rgb}{0.000000,0.000000,0.000000}%
\pgfsetstrokecolor{currentstroke}%
\pgfsetdash{}{0pt}%
\pgfsys@defobject{currentmarker}{\pgfqpoint{0.000000in}{-0.048611in}}{\pgfqpoint{0.000000in}{0.000000in}}{%
\pgfpathmoveto{\pgfqpoint{0.000000in}{0.000000in}}%
\pgfpathlineto{\pgfqpoint{0.000000in}{-0.048611in}}%
\pgfusepath{stroke,fill}%
}%
\begin{pgfscope}%
\pgfsys@transformshift{1.303236in}{0.450111in}%
\pgfsys@useobject{currentmarker}{}%
\end{pgfscope}%
\end{pgfscope}%
\begin{pgfscope}%
\definecolor{textcolor}{rgb}{0.000000,0.000000,0.000000}%
\pgfsetstrokecolor{textcolor}%
\pgfsetfillcolor{textcolor}%
\pgftext[x=1.303236in,y=0.352889in,,top]{\color{textcolor}{\rmfamily\fontsize{8.000000}{9.600000}\selectfont\catcode`\^=\active\def^{\ifmmode\sp\else\^{}\fi}\catcode`\%=\active\def%{\%}100}}%
\end{pgfscope}%
\begin{pgfscope}%
\pgfsetbuttcap%
\pgfsetroundjoin%
\definecolor{currentfill}{rgb}{0.000000,0.000000,0.000000}%
\pgfsetfillcolor{currentfill}%
\pgfsetlinewidth{0.803000pt}%
\definecolor{currentstroke}{rgb}{0.000000,0.000000,0.000000}%
\pgfsetstrokecolor{currentstroke}%
\pgfsetdash{}{0pt}%
\pgfsys@defobject{currentmarker}{\pgfqpoint{0.000000in}{-0.048611in}}{\pgfqpoint{0.000000in}{0.000000in}}{%
\pgfpathmoveto{\pgfqpoint{0.000000in}{0.000000in}}%
\pgfpathlineto{\pgfqpoint{0.000000in}{-0.048611in}}%
\pgfusepath{stroke,fill}%
}%
\begin{pgfscope}%
\pgfsys@transformshift{1.940983in}{0.450111in}%
\pgfsys@useobject{currentmarker}{}%
\end{pgfscope}%
\end{pgfscope}%
\begin{pgfscope}%
\definecolor{textcolor}{rgb}{0.000000,0.000000,0.000000}%
\pgfsetstrokecolor{textcolor}%
\pgfsetfillcolor{textcolor}%
\pgftext[x=1.940983in,y=0.352889in,,top]{\color{textcolor}{\rmfamily\fontsize{8.000000}{9.600000}\selectfont\catcode`\^=\active\def^{\ifmmode\sp\else\^{}\fi}\catcode`\%=\active\def%{\%}200}}%
\end{pgfscope}%
\begin{pgfscope}%
\pgfsetbuttcap%
\pgfsetroundjoin%
\definecolor{currentfill}{rgb}{0.000000,0.000000,0.000000}%
\pgfsetfillcolor{currentfill}%
\pgfsetlinewidth{0.803000pt}%
\definecolor{currentstroke}{rgb}{0.000000,0.000000,0.000000}%
\pgfsetstrokecolor{currentstroke}%
\pgfsetdash{}{0pt}%
\pgfsys@defobject{currentmarker}{\pgfqpoint{0.000000in}{-0.048611in}}{\pgfqpoint{0.000000in}{0.000000in}}{%
\pgfpathmoveto{\pgfqpoint{0.000000in}{0.000000in}}%
\pgfpathlineto{\pgfqpoint{0.000000in}{-0.048611in}}%
\pgfusepath{stroke,fill}%
}%
\begin{pgfscope}%
\pgfsys@transformshift{2.578730in}{0.450111in}%
\pgfsys@useobject{currentmarker}{}%
\end{pgfscope}%
\end{pgfscope}%
\begin{pgfscope}%
\definecolor{textcolor}{rgb}{0.000000,0.000000,0.000000}%
\pgfsetstrokecolor{textcolor}%
\pgfsetfillcolor{textcolor}%
\pgftext[x=2.578730in,y=0.352889in,,top]{\color{textcolor}{\rmfamily\fontsize{8.000000}{9.600000}\selectfont\catcode`\^=\active\def^{\ifmmode\sp\else\^{}\fi}\catcode`\%=\active\def%{\%}300}}%
\end{pgfscope}%
\begin{pgfscope}%
\pgfsetbuttcap%
\pgfsetroundjoin%
\definecolor{currentfill}{rgb}{0.000000,0.000000,0.000000}%
\pgfsetfillcolor{currentfill}%
\pgfsetlinewidth{0.803000pt}%
\definecolor{currentstroke}{rgb}{0.000000,0.000000,0.000000}%
\pgfsetstrokecolor{currentstroke}%
\pgfsetdash{}{0pt}%
\pgfsys@defobject{currentmarker}{\pgfqpoint{0.000000in}{-0.048611in}}{\pgfqpoint{0.000000in}{0.000000in}}{%
\pgfpathmoveto{\pgfqpoint{0.000000in}{0.000000in}}%
\pgfpathlineto{\pgfqpoint{0.000000in}{-0.048611in}}%
\pgfusepath{stroke,fill}%
}%
\begin{pgfscope}%
\pgfsys@transformshift{3.216477in}{0.450111in}%
\pgfsys@useobject{currentmarker}{}%
\end{pgfscope}%
\end{pgfscope}%
\begin{pgfscope}%
\definecolor{textcolor}{rgb}{0.000000,0.000000,0.000000}%
\pgfsetstrokecolor{textcolor}%
\pgfsetfillcolor{textcolor}%
\pgftext[x=3.216477in,y=0.352889in,,top]{\color{textcolor}{\rmfamily\fontsize{8.000000}{9.600000}\selectfont\catcode`\^=\active\def^{\ifmmode\sp\else\^{}\fi}\catcode`\%=\active\def%{\%}400}}%
\end{pgfscope}%
\begin{pgfscope}%
\pgfsetbuttcap%
\pgfsetroundjoin%
\definecolor{currentfill}{rgb}{0.000000,0.000000,0.000000}%
\pgfsetfillcolor{currentfill}%
\pgfsetlinewidth{0.803000pt}%
\definecolor{currentstroke}{rgb}{0.000000,0.000000,0.000000}%
\pgfsetstrokecolor{currentstroke}%
\pgfsetdash{}{0pt}%
\pgfsys@defobject{currentmarker}{\pgfqpoint{0.000000in}{-0.048611in}}{\pgfqpoint{0.000000in}{0.000000in}}{%
\pgfpathmoveto{\pgfqpoint{0.000000in}{0.000000in}}%
\pgfpathlineto{\pgfqpoint{0.000000in}{-0.048611in}}%
\pgfusepath{stroke,fill}%
}%
\begin{pgfscope}%
\pgfsys@transformshift{3.854223in}{0.450111in}%
\pgfsys@useobject{currentmarker}{}%
\end{pgfscope}%
\end{pgfscope}%
\begin{pgfscope}%
\definecolor{textcolor}{rgb}{0.000000,0.000000,0.000000}%
\pgfsetstrokecolor{textcolor}%
\pgfsetfillcolor{textcolor}%
\pgftext[x=3.854223in,y=0.352889in,,top]{\color{textcolor}{\rmfamily\fontsize{8.000000}{9.600000}\selectfont\catcode`\^=\active\def^{\ifmmode\sp\else\^{}\fi}\catcode`\%=\active\def%{\%}500}}%
\end{pgfscope}%
\begin{pgfscope}%
\pgfsetbuttcap%
\pgfsetroundjoin%
\definecolor{currentfill}{rgb}{0.000000,0.000000,0.000000}%
\pgfsetfillcolor{currentfill}%
\pgfsetlinewidth{0.803000pt}%
\definecolor{currentstroke}{rgb}{0.000000,0.000000,0.000000}%
\pgfsetstrokecolor{currentstroke}%
\pgfsetdash{}{0pt}%
\pgfsys@defobject{currentmarker}{\pgfqpoint{0.000000in}{-0.048611in}}{\pgfqpoint{0.000000in}{0.000000in}}{%
\pgfpathmoveto{\pgfqpoint{0.000000in}{0.000000in}}%
\pgfpathlineto{\pgfqpoint{0.000000in}{-0.048611in}}%
\pgfusepath{stroke,fill}%
}%
\begin{pgfscope}%
\pgfsys@transformshift{4.491970in}{0.450111in}%
\pgfsys@useobject{currentmarker}{}%
\end{pgfscope}%
\end{pgfscope}%
\begin{pgfscope}%
\definecolor{textcolor}{rgb}{0.000000,0.000000,0.000000}%
\pgfsetstrokecolor{textcolor}%
\pgfsetfillcolor{textcolor}%
\pgftext[x=4.491970in,y=0.352889in,,top]{\color{textcolor}{\rmfamily\fontsize{8.000000}{9.600000}\selectfont\catcode`\^=\active\def^{\ifmmode\sp\else\^{}\fi}\catcode`\%=\active\def%{\%}600}}%
\end{pgfscope}%
\begin{pgfscope}%
\pgfsetbuttcap%
\pgfsetroundjoin%
\definecolor{currentfill}{rgb}{0.000000,0.000000,0.000000}%
\pgfsetfillcolor{currentfill}%
\pgfsetlinewidth{0.803000pt}%
\definecolor{currentstroke}{rgb}{0.000000,0.000000,0.000000}%
\pgfsetstrokecolor{currentstroke}%
\pgfsetdash{}{0pt}%
\pgfsys@defobject{currentmarker}{\pgfqpoint{0.000000in}{-0.048611in}}{\pgfqpoint{0.000000in}{0.000000in}}{%
\pgfpathmoveto{\pgfqpoint{0.000000in}{0.000000in}}%
\pgfpathlineto{\pgfqpoint{0.000000in}{-0.048611in}}%
\pgfusepath{stroke,fill}%
}%
\begin{pgfscope}%
\pgfsys@transformshift{5.129717in}{0.450111in}%
\pgfsys@useobject{currentmarker}{}%
\end{pgfscope}%
\end{pgfscope}%
\begin{pgfscope}%
\definecolor{textcolor}{rgb}{0.000000,0.000000,0.000000}%
\pgfsetstrokecolor{textcolor}%
\pgfsetfillcolor{textcolor}%
\pgftext[x=5.129717in,y=0.352889in,,top]{\color{textcolor}{\rmfamily\fontsize{8.000000}{9.600000}\selectfont\catcode`\^=\active\def^{\ifmmode\sp\else\^{}\fi}\catcode`\%=\active\def%{\%}700}}%
\end{pgfscope}%
\begin{pgfscope}%
\definecolor{textcolor}{rgb}{0.000000,0.000000,0.000000}%
\pgfsetstrokecolor{textcolor}%
\pgfsetfillcolor{textcolor}%
\pgftext[x=2.999643in,y=0.198667in,,top]{\color{textcolor}{\rmfamily\fontsize{8.000000}{9.600000}\selectfont\catcode`\^=\active\def^{\ifmmode\sp\else\^{}\fi}\catcode`\%=\active\def%{\%}Time $t$ in d}}%
\end{pgfscope}%
\begin{pgfscope}%
\pgfsetbuttcap%
\pgfsetroundjoin%
\definecolor{currentfill}{rgb}{0.000000,0.000000,0.000000}%
\pgfsetfillcolor{currentfill}%
\pgfsetlinewidth{0.803000pt}%
\definecolor{currentstroke}{rgb}{0.000000,0.000000,0.000000}%
\pgfsetstrokecolor{currentstroke}%
\pgfsetdash{}{0pt}%
\pgfsys@defobject{currentmarker}{\pgfqpoint{-0.048611in}{0.000000in}}{\pgfqpoint{-0.000000in}{0.000000in}}{%
\pgfpathmoveto{\pgfqpoint{-0.000000in}{0.000000in}}%
\pgfpathlineto{\pgfqpoint{-0.048611in}{0.000000in}}%
\pgfusepath{stroke,fill}%
}%
\begin{pgfscope}%
\pgfsys@transformshift{0.665489in}{0.546740in}%
\pgfsys@useobject{currentmarker}{}%
\end{pgfscope}%
\end{pgfscope}%
\begin{pgfscope}%
\definecolor{textcolor}{rgb}{0.000000,0.000000,0.000000}%
\pgfsetstrokecolor{textcolor}%
\pgfsetfillcolor{textcolor}%
\pgftext[x=0.266667in, y=0.508184in, left, base]{\color{textcolor}{\rmfamily\fontsize{8.000000}{9.600000}\selectfont\catcode`\^=\active\def^{\ifmmode\sp\else\^{}\fi}\catcode`\%=\active\def%{\%}\ensuremath{-}0.50}}%
\end{pgfscope}%
\begin{pgfscope}%
\pgfsetbuttcap%
\pgfsetroundjoin%
\definecolor{currentfill}{rgb}{0.000000,0.000000,0.000000}%
\pgfsetfillcolor{currentfill}%
\pgfsetlinewidth{0.803000pt}%
\definecolor{currentstroke}{rgb}{0.000000,0.000000,0.000000}%
\pgfsetstrokecolor{currentstroke}%
\pgfsetdash{}{0pt}%
\pgfsys@defobject{currentmarker}{\pgfqpoint{-0.048611in}{0.000000in}}{\pgfqpoint{-0.000000in}{0.000000in}}{%
\pgfpathmoveto{\pgfqpoint{-0.000000in}{0.000000in}}%
\pgfpathlineto{\pgfqpoint{-0.048611in}{0.000000in}}%
\pgfusepath{stroke,fill}%
}%
\begin{pgfscope}%
\pgfsys@transformshift{0.665489in}{0.788312in}%
\pgfsys@useobject{currentmarker}{}%
\end{pgfscope}%
\end{pgfscope}%
\begin{pgfscope}%
\definecolor{textcolor}{rgb}{0.000000,0.000000,0.000000}%
\pgfsetstrokecolor{textcolor}%
\pgfsetfillcolor{textcolor}%
\pgftext[x=0.266667in, y=0.749757in, left, base]{\color{textcolor}{\rmfamily\fontsize{8.000000}{9.600000}\selectfont\catcode`\^=\active\def^{\ifmmode\sp\else\^{}\fi}\catcode`\%=\active\def%{\%}\ensuremath{-}0.25}}%
\end{pgfscope}%
\begin{pgfscope}%
\pgfsetbuttcap%
\pgfsetroundjoin%
\definecolor{currentfill}{rgb}{0.000000,0.000000,0.000000}%
\pgfsetfillcolor{currentfill}%
\pgfsetlinewidth{0.803000pt}%
\definecolor{currentstroke}{rgb}{0.000000,0.000000,0.000000}%
\pgfsetstrokecolor{currentstroke}%
\pgfsetdash{}{0pt}%
\pgfsys@defobject{currentmarker}{\pgfqpoint{-0.048611in}{0.000000in}}{\pgfqpoint{-0.000000in}{0.000000in}}{%
\pgfpathmoveto{\pgfqpoint{-0.000000in}{0.000000in}}%
\pgfpathlineto{\pgfqpoint{-0.048611in}{0.000000in}}%
\pgfusepath{stroke,fill}%
}%
\begin{pgfscope}%
\pgfsys@transformshift{0.665489in}{1.029884in}%
\pgfsys@useobject{currentmarker}{}%
\end{pgfscope}%
\end{pgfscope}%
\begin{pgfscope}%
\definecolor{textcolor}{rgb}{0.000000,0.000000,0.000000}%
\pgfsetstrokecolor{textcolor}%
\pgfsetfillcolor{textcolor}%
\pgftext[x=0.358489in, y=0.991329in, left, base]{\color{textcolor}{\rmfamily\fontsize{8.000000}{9.600000}\selectfont\catcode`\^=\active\def^{\ifmmode\sp\else\^{}\fi}\catcode`\%=\active\def%{\%}0.00}}%
\end{pgfscope}%
\begin{pgfscope}%
\pgfsetbuttcap%
\pgfsetroundjoin%
\definecolor{currentfill}{rgb}{0.000000,0.000000,0.000000}%
\pgfsetfillcolor{currentfill}%
\pgfsetlinewidth{0.803000pt}%
\definecolor{currentstroke}{rgb}{0.000000,0.000000,0.000000}%
\pgfsetstrokecolor{currentstroke}%
\pgfsetdash{}{0pt}%
\pgfsys@defobject{currentmarker}{\pgfqpoint{-0.048611in}{0.000000in}}{\pgfqpoint{-0.000000in}{0.000000in}}{%
\pgfpathmoveto{\pgfqpoint{-0.000000in}{0.000000in}}%
\pgfpathlineto{\pgfqpoint{-0.048611in}{0.000000in}}%
\pgfusepath{stroke,fill}%
}%
\begin{pgfscope}%
\pgfsys@transformshift{0.665489in}{1.271457in}%
\pgfsys@useobject{currentmarker}{}%
\end{pgfscope}%
\end{pgfscope}%
\begin{pgfscope}%
\definecolor{textcolor}{rgb}{0.000000,0.000000,0.000000}%
\pgfsetstrokecolor{textcolor}%
\pgfsetfillcolor{textcolor}%
\pgftext[x=0.358489in, y=1.232901in, left, base]{\color{textcolor}{\rmfamily\fontsize{8.000000}{9.600000}\selectfont\catcode`\^=\active\def^{\ifmmode\sp\else\^{}\fi}\catcode`\%=\active\def%{\%}0.25}}%
\end{pgfscope}%
\begin{pgfscope}%
\pgfsetbuttcap%
\pgfsetroundjoin%
\definecolor{currentfill}{rgb}{0.000000,0.000000,0.000000}%
\pgfsetfillcolor{currentfill}%
\pgfsetlinewidth{0.803000pt}%
\definecolor{currentstroke}{rgb}{0.000000,0.000000,0.000000}%
\pgfsetstrokecolor{currentstroke}%
\pgfsetdash{}{0pt}%
\pgfsys@defobject{currentmarker}{\pgfqpoint{-0.048611in}{0.000000in}}{\pgfqpoint{-0.000000in}{0.000000in}}{%
\pgfpathmoveto{\pgfqpoint{-0.000000in}{0.000000in}}%
\pgfpathlineto{\pgfqpoint{-0.048611in}{0.000000in}}%
\pgfusepath{stroke,fill}%
}%
\begin{pgfscope}%
\pgfsys@transformshift{0.665489in}{1.513029in}%
\pgfsys@useobject{currentmarker}{}%
\end{pgfscope}%
\end{pgfscope}%
\begin{pgfscope}%
\definecolor{textcolor}{rgb}{0.000000,0.000000,0.000000}%
\pgfsetstrokecolor{textcolor}%
\pgfsetfillcolor{textcolor}%
\pgftext[x=0.358489in, y=1.474474in, left, base]{\color{textcolor}{\rmfamily\fontsize{8.000000}{9.600000}\selectfont\catcode`\^=\active\def^{\ifmmode\sp\else\^{}\fi}\catcode`\%=\active\def%{\%}0.50}}%
\end{pgfscope}%
\begin{pgfscope}%
\definecolor{textcolor}{rgb}{0.000000,0.000000,0.000000}%
\pgfsetstrokecolor{textcolor}%
\pgfsetfillcolor{textcolor}%
\pgftext[x=0.211111in,y=1.029884in,,bottom,rotate=90.000000]{\color{textcolor}{\rmfamily\fontsize{8.000000}{9.600000}\selectfont\catcode`\^=\active\def^{\ifmmode\sp\else\^{}\fi}\catcode`\%=\active\def%{\%}$c(t)$}}%
\end{pgfscope}%
\begin{pgfscope}%
\pgfpathrectangle{\pgfqpoint{0.665489in}{0.450111in}}{\pgfqpoint{4.668307in}{1.159547in}}%
\pgfusepath{clip}%
\pgfsetrectcap%
\pgfsetroundjoin%
\pgfsetlinewidth{0.501875pt}%
\definecolor{currentstroke}{rgb}{0.301961,0.301961,0.301961}%
\pgfsetstrokecolor{currentstroke}%
\pgfsetdash{}{0pt}%
\pgfpathmoveto{\pgfqpoint{0.665489in}{1.029884in}}%
\pgfpathlineto{\pgfqpoint{5.333796in}{1.029884in}}%
\pgfusepath{stroke}%
\end{pgfscope}%
\begin{pgfscope}%
\pgfpathrectangle{\pgfqpoint{0.665489in}{0.450111in}}{\pgfqpoint{4.668307in}{1.159547in}}%
\pgfusepath{clip}%
\pgfsetrectcap%
\pgfsetroundjoin%
\pgfsetlinewidth{1.003750pt}%
\definecolor{currentstroke}{rgb}{0.007843,0.619608,0.450980}%
\pgfsetstrokecolor{currentstroke}%
\pgfsetdash{}{0pt}%
\pgfpathmoveto{\pgfqpoint{0.665622in}{1.502929in}}%
\pgfpathlineto{\pgfqpoint{0.677978in}{1.478135in}}%
\pgfpathlineto{\pgfqpoint{0.702159in}{1.430031in}}%
\pgfpathlineto{\pgfqpoint{0.705481in}{1.423367in}}%
\pgfpathlineto{\pgfqpoint{0.714649in}{1.405034in}}%
\pgfpathlineto{\pgfqpoint{0.718369in}{1.397703in}}%
\pgfpathlineto{\pgfqpoint{0.722089in}{1.390106in}}%
\pgfpathlineto{\pgfqpoint{0.724481in}{1.385433in}}%
\pgfpathlineto{\pgfqpoint{0.735907in}{1.362485in}}%
\pgfpathlineto{\pgfqpoint{0.737900in}{1.358579in}}%
\pgfpathlineto{\pgfqpoint{0.749326in}{1.335774in}}%
\pgfpathlineto{\pgfqpoint{0.771382in}{1.291707in}}%
\pgfpathlineto{\pgfqpoint{0.775633in}{1.283206in}}%
\pgfpathlineto{\pgfqpoint{0.783472in}{1.267515in}}%
\pgfpathlineto{\pgfqpoint{0.790647in}{1.253251in}}%
\pgfpathlineto{\pgfqpoint{0.803668in}{1.227166in}}%
\pgfpathlineto{\pgfqpoint{0.808318in}{1.217884in}}%
\pgfpathlineto{\pgfqpoint{0.823199in}{1.188190in}}%
\pgfpathlineto{\pgfqpoint{0.834359in}{1.165889in}}%
\pgfpathlineto{\pgfqpoint{0.928427in}{0.977967in}}%
\pgfpathlineto{\pgfqpoint{0.956992in}{0.920901in}}%
\pgfpathlineto{\pgfqpoint{1.115499in}{0.604156in}}%
\pgfpathmoveto{\pgfqpoint{1.115632in}{1.277601in}}%
\pgfpathlineto{\pgfqpoint{1.119751in}{1.272374in}}%
\pgfpathlineto{\pgfqpoint{1.122142in}{1.269399in}}%
\pgfpathlineto{\pgfqpoint{1.127457in}{1.262602in}}%
\pgfpathlineto{\pgfqpoint{1.130380in}{1.258968in}}%
\pgfpathlineto{\pgfqpoint{1.135562in}{1.252372in}}%
\pgfpathlineto{\pgfqpoint{1.140610in}{1.245981in}}%
\pgfpathlineto{\pgfqpoint{1.145925in}{1.239189in}}%
\pgfpathlineto{\pgfqpoint{1.148848in}{1.235519in}}%
\pgfpathlineto{\pgfqpoint{1.183260in}{1.192016in}}%
\pgfpathlineto{\pgfqpoint{1.191896in}{1.181079in}}%
\pgfpathlineto{\pgfqpoint{1.214616in}{1.152333in}}%
\pgfpathlineto{\pgfqpoint{1.222189in}{1.142757in}}%
\pgfpathlineto{\pgfqpoint{1.237335in}{1.123585in}}%
\pgfpathlineto{\pgfqpoint{1.267894in}{1.084929in}}%
\pgfpathlineto{\pgfqpoint{1.350535in}{0.980358in}}%
\pgfpathlineto{\pgfqpoint{1.365682in}{0.961223in}}%
\pgfpathlineto{\pgfqpoint{1.380031in}{0.943046in}}%
\pgfpathlineto{\pgfqpoint{1.407401in}{0.908407in}}%
\pgfpathlineto{\pgfqpoint{1.411254in}{0.903573in}}%
\pgfpathlineto{\pgfqpoint{1.417233in}{0.896012in}}%
\pgfpathlineto{\pgfqpoint{1.536412in}{0.745212in}}%
\pgfpathlineto{\pgfqpoint{1.540664in}{0.739853in}}%
\pgfpathlineto{\pgfqpoint{1.580789in}{0.689009in}}%
\pgfpathlineto{\pgfqpoint{1.581719in}{0.687917in}}%
\pgfpathlineto{\pgfqpoint{1.586900in}{0.681337in}}%
\pgfpathlineto{\pgfqpoint{1.591816in}{0.675122in}}%
\pgfpathlineto{\pgfqpoint{1.595271in}{0.670681in}}%
\pgfpathlineto{\pgfqpoint{1.598061in}{0.667177in}}%
\pgfpathlineto{\pgfqpoint{1.636060in}{0.619122in}}%
\pgfpathmoveto{\pgfqpoint{1.636193in}{1.260765in}}%
\pgfpathlineto{\pgfqpoint{1.641242in}{1.253342in}}%
\pgfpathlineto{\pgfqpoint{1.642969in}{1.250903in}}%
\pgfpathlineto{\pgfqpoint{1.652934in}{1.236272in}}%
\pgfpathlineto{\pgfqpoint{1.654130in}{1.234625in}}%
\pgfpathlineto{\pgfqpoint{1.665689in}{1.217692in}}%
\pgfpathlineto{\pgfqpoint{1.675786in}{1.202958in}}%
\pgfpathlineto{\pgfqpoint{1.876411in}{0.910070in}}%
\pgfpathlineto{\pgfqpoint{1.894215in}{0.884064in}}%
\pgfpathlineto{\pgfqpoint{1.898333in}{0.878092in}}%
\pgfpathlineto{\pgfqpoint{1.904445in}{0.869162in}}%
\pgfpathlineto{\pgfqpoint{1.916802in}{0.851122in}}%
\pgfpathlineto{\pgfqpoint{1.922116in}{0.843328in}}%
\pgfpathlineto{\pgfqpoint{1.930619in}{0.830947in}}%
\pgfpathlineto{\pgfqpoint{1.942179in}{0.814075in}}%
\pgfpathlineto{\pgfqpoint{1.954933in}{0.795487in}}%
\pgfpathlineto{\pgfqpoint{1.959584in}{0.788670in}}%
\pgfpathlineto{\pgfqpoint{1.964101in}{0.782092in}}%
\pgfpathlineto{\pgfqpoint{1.971807in}{0.770770in}}%
\pgfpathlineto{\pgfqpoint{1.974199in}{0.767285in}}%
\pgfpathlineto{\pgfqpoint{1.980178in}{0.758603in}}%
\pgfpathlineto{\pgfqpoint{1.982835in}{0.754727in}}%
\pgfpathlineto{\pgfqpoint{1.985625in}{0.750550in}}%
\pgfpathlineto{\pgfqpoint{1.987352in}{0.748136in}}%
\pgfpathlineto{\pgfqpoint{1.992268in}{0.740877in}}%
\pgfpathlineto{\pgfqpoint{1.993464in}{0.739274in}}%
\pgfpathlineto{\pgfqpoint{1.997317in}{0.733593in}}%
\pgfpathlineto{\pgfqpoint{2.008212in}{0.717615in}}%
\pgfpathlineto{\pgfqpoint{2.010869in}{0.713945in}}%
\pgfpathlineto{\pgfqpoint{2.014058in}{0.709079in}}%
\pgfpathlineto{\pgfqpoint{2.019372in}{0.701532in}}%
\pgfpathlineto{\pgfqpoint{2.021498in}{0.698456in}}%
\pgfpathlineto{\pgfqpoint{2.024023in}{0.694498in}}%
\pgfpathlineto{\pgfqpoint{2.024687in}{0.693645in}}%
\pgfpathlineto{\pgfqpoint{2.034918in}{0.678793in}}%
\pgfpathlineto{\pgfqpoint{2.043952in}{0.665507in}}%
\pgfpathlineto{\pgfqpoint{2.046875in}{0.661214in}}%
\pgfpathlineto{\pgfqpoint{2.050064in}{0.656593in}}%
\pgfpathlineto{\pgfqpoint{2.055113in}{0.649150in}}%
\pgfpathlineto{\pgfqpoint{2.057106in}{0.646237in}}%
\pgfpathlineto{\pgfqpoint{2.060029in}{0.641980in}}%
\pgfpathlineto{\pgfqpoint{2.061623in}{0.639689in}}%
\pgfpathlineto{\pgfqpoint{2.067204in}{0.631511in}}%
\pgfpathlineto{\pgfqpoint{2.070525in}{0.626708in}}%
\pgfpathlineto{\pgfqpoint{2.079826in}{0.613171in}}%
\pgfpathlineto{\pgfqpoint{2.082350in}{0.609496in}}%
\pgfpathlineto{\pgfqpoint{2.085140in}{0.605347in}}%
\pgfpathlineto{\pgfqpoint{2.087665in}{0.601672in}}%
\pgfpathlineto{\pgfqpoint{2.097762in}{0.586943in}}%
\pgfpathmoveto{\pgfqpoint{2.097895in}{1.233733in}}%
\pgfpathlineto{\pgfqpoint{2.104538in}{1.224597in}}%
\pgfpathlineto{\pgfqpoint{2.116895in}{1.207712in}}%
\pgfpathlineto{\pgfqpoint{2.126727in}{1.194216in}}%
\pgfpathlineto{\pgfqpoint{2.130181in}{1.189494in}}%
\pgfpathlineto{\pgfqpoint{2.165124in}{1.141626in}}%
\pgfpathlineto{\pgfqpoint{2.197676in}{1.097043in}}%
\pgfpathlineto{\pgfqpoint{2.218934in}{1.067888in}}%
\pgfpathlineto{\pgfqpoint{2.236738in}{1.043526in}}%
\pgfpathlineto{\pgfqpoint{2.320442in}{0.928874in}}%
\pgfpathlineto{\pgfqpoint{2.329078in}{0.917006in}}%
\pgfpathlineto{\pgfqpoint{2.473236in}{0.719522in}}%
\pgfpathlineto{\pgfqpoint{2.474697in}{0.717554in}}%
\pgfpathlineto{\pgfqpoint{2.479746in}{0.710599in}}%
\pgfpathlineto{\pgfqpoint{2.496088in}{0.688202in}}%
\pgfpathlineto{\pgfqpoint{2.498480in}{0.684973in}}%
\pgfpathlineto{\pgfqpoint{2.501934in}{0.680216in}}%
\pgfpathlineto{\pgfqpoint{2.512962in}{0.665112in}}%
\pgfpathlineto{\pgfqpoint{2.517081in}{0.659479in}}%
\pgfpathlineto{\pgfqpoint{2.522395in}{0.652228in}}%
\pgfpathlineto{\pgfqpoint{2.525717in}{0.647596in}}%
\pgfpathlineto{\pgfqpoint{2.527178in}{0.645677in}}%
\pgfpathlineto{\pgfqpoint{2.530367in}{0.641183in}}%
\pgfpathmoveto{\pgfqpoint{2.530500in}{1.267602in}}%
\pgfpathlineto{\pgfqpoint{2.615001in}{1.142235in}}%
\pgfpathlineto{\pgfqpoint{2.624435in}{1.128228in}}%
\pgfpathlineto{\pgfqpoint{2.642770in}{1.101050in}}%
\pgfpathlineto{\pgfqpoint{2.704419in}{1.009597in}}%
\pgfpathlineto{\pgfqpoint{2.727139in}{0.975913in}}%
\pgfpathlineto{\pgfqpoint{2.741488in}{0.954618in}}%
\pgfpathlineto{\pgfqpoint{2.911554in}{0.702405in}}%
\pgfpathlineto{\pgfqpoint{2.914743in}{0.697615in}}%
\pgfpathlineto{\pgfqpoint{2.926567in}{0.680072in}}%
\pgfpathlineto{\pgfqpoint{2.929358in}{0.675911in}}%
\pgfpathlineto{\pgfqpoint{2.935602in}{0.666731in}}%
\pgfpathlineto{\pgfqpoint{2.939455in}{0.660988in}}%
\pgfpathlineto{\pgfqpoint{2.949420in}{0.646174in}}%
\pgfpathlineto{\pgfqpoint{2.956728in}{0.635399in}}%
\pgfpathmoveto{\pgfqpoint{2.956860in}{1.208149in}}%
\pgfpathlineto{\pgfqpoint{2.966294in}{1.201097in}}%
\pgfpathlineto{\pgfqpoint{2.971874in}{1.196764in}}%
\pgfpathlineto{\pgfqpoint{2.973070in}{1.195929in}}%
\pgfpathlineto{\pgfqpoint{2.980510in}{1.190320in}}%
\pgfpathlineto{\pgfqpoint{3.034054in}{1.149872in}}%
\pgfpathlineto{\pgfqpoint{3.041760in}{1.144054in}}%
\pgfpathlineto{\pgfqpoint{3.048404in}{1.139017in}}%
\pgfpathlineto{\pgfqpoint{3.053320in}{1.135322in}}%
\pgfpathlineto{\pgfqpoint{3.072851in}{1.120580in}}%
\pgfpathlineto{\pgfqpoint{3.120017in}{1.084912in}}%
\pgfpathlineto{\pgfqpoint{3.135695in}{1.073062in}}%
\pgfpathlineto{\pgfqpoint{3.177282in}{1.041646in}}%
\pgfpathlineto{\pgfqpoint{3.409395in}{0.866277in}}%
\pgfpathlineto{\pgfqpoint{3.412318in}{0.864061in}}%
\pgfpathlineto{\pgfqpoint{3.416437in}{0.860898in}}%
\pgfpathlineto{\pgfqpoint{3.418031in}{0.859721in}}%
\pgfpathlineto{\pgfqpoint{3.428527in}{0.851798in}}%
\pgfpathlineto{\pgfqpoint{3.433842in}{0.847811in}}%
\pgfpathlineto{\pgfqpoint{3.436765in}{0.845360in}}%
\pgfpathlineto{\pgfqpoint{3.437828in}{0.844817in}}%
\pgfpathlineto{\pgfqpoint{3.442877in}{0.840981in}}%
\pgfpathlineto{\pgfqpoint{3.445135in}{0.839295in}}%
\pgfpathlineto{\pgfqpoint{3.453904in}{0.832622in}}%
\pgfpathlineto{\pgfqpoint{3.457226in}{0.830205in}}%
\pgfpathlineto{\pgfqpoint{3.459485in}{0.828369in}}%
\pgfpathlineto{\pgfqpoint{3.462275in}{0.826282in}}%
\pgfpathlineto{\pgfqpoint{3.477156in}{0.815099in}}%
\pgfpathmoveto{\pgfqpoint{3.477288in}{1.325044in}}%
\pgfpathlineto{\pgfqpoint{3.541063in}{1.242309in}}%
\pgfpathlineto{\pgfqpoint{3.543720in}{1.238904in}}%
\pgfpathlineto{\pgfqpoint{3.553951in}{1.225591in}}%
\pgfpathlineto{\pgfqpoint{3.560461in}{1.217198in}}%
\pgfpathlineto{\pgfqpoint{3.567636in}{1.207852in}}%
\pgfpathlineto{\pgfqpoint{3.671934in}{1.072531in}}%
\pgfpathlineto{\pgfqpoint{3.902852in}{0.772909in}}%
\pgfpathlineto{\pgfqpoint{3.903649in}{0.771893in}}%
\pgfpathlineto{\pgfqpoint{3.911621in}{0.761563in}}%
\pgfpathmoveto{\pgfqpoint{3.911753in}{1.414503in}}%
\pgfpathlineto{\pgfqpoint{3.914145in}{1.411445in}}%
\pgfpathlineto{\pgfqpoint{3.915739in}{1.409488in}}%
\pgfpathlineto{\pgfqpoint{3.924243in}{1.398820in}}%
\pgfpathlineto{\pgfqpoint{3.926236in}{1.396365in}}%
\pgfpathlineto{\pgfqpoint{4.119686in}{1.154380in}}%
\pgfpathlineto{\pgfqpoint{4.122609in}{1.150696in}}%
\pgfpathlineto{\pgfqpoint{4.132175in}{1.138721in}}%
\pgfpathlineto{\pgfqpoint{4.177216in}{1.082400in}}%
\pgfpathlineto{\pgfqpoint{4.190502in}{1.065761in}}%
\pgfpathlineto{\pgfqpoint{4.214285in}{1.036015in}}%
\pgfpathlineto{\pgfqpoint{4.386343in}{0.820741in}}%
\pgfpathmoveto{\pgfqpoint{4.386476in}{1.254392in}}%
\pgfpathlineto{\pgfqpoint{4.394182in}{1.244784in}}%
\pgfpathlineto{\pgfqpoint{4.398700in}{1.239165in}}%
\pgfpathlineto{\pgfqpoint{4.405476in}{1.230680in}}%
\pgfpathlineto{\pgfqpoint{4.408000in}{1.227556in}}%
\pgfpathlineto{\pgfqpoint{4.415573in}{1.218068in}}%
\pgfpathlineto{\pgfqpoint{4.420622in}{1.211804in}}%
\pgfpathlineto{\pgfqpoint{4.438958in}{1.188929in}}%
\pgfpathlineto{\pgfqpoint{4.546577in}{1.054663in}}%
\pgfpathlineto{\pgfqpoint{4.566640in}{1.029625in}}%
\pgfpathlineto{\pgfqpoint{4.814431in}{0.720413in}}%
\pgfpathmoveto{\pgfqpoint{4.814564in}{1.352811in}}%
\pgfpathlineto{\pgfqpoint{4.832235in}{1.326179in}}%
\pgfpathlineto{\pgfqpoint{4.834361in}{1.323027in}}%
\pgfpathlineto{\pgfqpoint{4.839409in}{1.315401in}}%
\pgfpathlineto{\pgfqpoint{4.852696in}{1.295468in}}%
\pgfpathlineto{\pgfqpoint{4.923512in}{1.189111in}}%
\pgfpathlineto{\pgfqpoint{4.927365in}{1.183329in}}%
\pgfpathlineto{\pgfqpoint{4.970546in}{1.118370in}}%
\pgfpathlineto{\pgfqpoint{4.978917in}{1.105832in}}%
\pgfpathlineto{\pgfqpoint{4.987686in}{1.092638in}}%
\pgfpathlineto{\pgfqpoint{5.000175in}{1.073870in}}%
\pgfpathlineto{\pgfqpoint{5.015720in}{1.050508in}}%
\pgfpathlineto{\pgfqpoint{5.082949in}{0.949467in}}%
\pgfpathlineto{\pgfqpoint{5.091718in}{0.936314in}}%
\pgfpathlineto{\pgfqpoint{5.102746in}{0.919709in}}%
\pgfpathlineto{\pgfqpoint{5.118291in}{0.896402in}}%
\pgfpathlineto{\pgfqpoint{5.123871in}{0.887991in}}%
\pgfpathlineto{\pgfqpoint{5.154297in}{0.842280in}}%
\pgfpathlineto{\pgfqpoint{5.160542in}{0.832865in}}%
\pgfpathlineto{\pgfqpoint{5.167849in}{0.821949in}}%
\pgfpathlineto{\pgfqpoint{5.181534in}{0.801356in}}%
\pgfpathlineto{\pgfqpoint{5.210498in}{0.757753in}}%
\pgfpathlineto{\pgfqpoint{5.211827in}{0.755709in}}%
\pgfpathlineto{\pgfqpoint{5.233218in}{0.723656in}}%
\pgfpathlineto{\pgfqpoint{5.260455in}{0.682680in}}%
\pgfpathlineto{\pgfqpoint{5.261385in}{0.681348in}}%
\pgfpathlineto{\pgfqpoint{5.278658in}{0.655380in}}%
\pgfpathlineto{\pgfqpoint{5.280783in}{0.652030in}}%
\pgfpathlineto{\pgfqpoint{5.280916in}{0.652083in}}%
\pgfpathlineto{\pgfqpoint{5.284371in}{0.646766in}}%
\pgfpathlineto{\pgfqpoint{5.292741in}{0.634206in}}%
\pgfpathmoveto{\pgfqpoint{5.292874in}{1.265935in}}%
\pgfpathlineto{\pgfqpoint{5.298454in}{1.260716in}}%
\pgfpathlineto{\pgfqpoint{5.300182in}{1.259218in}}%
\pgfpathlineto{\pgfqpoint{5.305230in}{1.254474in}}%
\pgfpathlineto{\pgfqpoint{5.306825in}{1.253057in}}%
\pgfpathlineto{\pgfqpoint{5.313999in}{1.246520in}}%
\pgfpathlineto{\pgfqpoint{5.317188in}{1.243564in}}%
\pgfpathlineto{\pgfqpoint{5.322636in}{1.238570in}}%
\pgfpathlineto{\pgfqpoint{5.327817in}{1.233724in}}%
\pgfpathlineto{\pgfqpoint{5.330873in}{1.230959in}}%
\pgfpathlineto{\pgfqpoint{5.333796in}{1.228308in}}%
\pgfpathlineto{\pgfqpoint{5.333796in}{1.228308in}}%
\pgfusepath{stroke}%
\end{pgfscope}%
\begin{pgfscope}%
\pgfsetrectcap%
\pgfsetmiterjoin%
\pgfsetlinewidth{0.803000pt}%
\definecolor{currentstroke}{rgb}{0.000000,0.000000,0.000000}%
\pgfsetstrokecolor{currentstroke}%
\pgfsetdash{}{0pt}%
\pgfpathmoveto{\pgfqpoint{0.665489in}{0.450111in}}%
\pgfpathlineto{\pgfqpoint{0.665489in}{1.609658in}}%
\pgfusepath{stroke}%
\end{pgfscope}%
\begin{pgfscope}%
\pgfsetrectcap%
\pgfsetmiterjoin%
\pgfsetlinewidth{0.803000pt}%
\definecolor{currentstroke}{rgb}{0.000000,0.000000,0.000000}%
\pgfsetstrokecolor{currentstroke}%
\pgfsetdash{}{0pt}%
\pgfpathmoveto{\pgfqpoint{0.665489in}{0.450111in}}%
\pgfpathlineto{\pgfqpoint{5.333796in}{0.450111in}}%
\pgfusepath{stroke}%
\end{pgfscope}%
\begin{pgfscope}%
\definecolor{textcolor}{rgb}{0.000000,0.000000,0.000000}%
\pgfsetstrokecolor{textcolor}%
\pgfsetfillcolor{textcolor}%
\pgftext[x=5.333796in,y=1.331850in,left,]{\color{textcolor}{\rmfamily\fontsize{8.000000}{9.600000}\selectfont\catcode`\^=\active\def^{\ifmmode\sp\else\^{}\fi}\catcode`\%=\active\def%{\%} upper band}}%
\end{pgfscope}%
\begin{pgfscope}%
\definecolor{textcolor}{rgb}{0.000000,0.000000,0.000000}%
\pgfsetstrokecolor{textcolor}%
\pgfsetfillcolor{textcolor}%
\pgftext[x=5.333796in,y=0.727919in,left,]{\color{textcolor}{\rmfamily\fontsize{8.000000}{9.600000}\selectfont\catcode`\^=\active\def^{\ifmmode\sp\else\^{}\fi}\catcode`\%=\active\def%{\%} lower band}}%
\end{pgfscope}%
\end{pgfpicture}%
\makeatother%
\endgroup%

    \caption{Drift signal $c(t)$ of \texttt{IDV(29)}.}
    \label{fig:drift_signal}
\end{figure}
%
Table~\ref{tab:idv29} lists the parameters of \texttt{IDV(29)}. The segment duration of $\SI{365}{\day}/5$ represents five maintenance events per year. This results in a relatively large number of operating points per calibration interval, so that the calibration error changes only slightly when an operating point is reached. Consequently, each operating point reflects a single drift value. The shift fraction keeps adjacent segment boundaries apart while still breaking the periodicity of the grid. The drift coefficient is bounded from above by the requirement that the peak of the calibration error remains below the spread of the stripper temperature caused by the process itself. With the chosen value, the recorded temperature deviates from the process value by at most \idvPeakErrorPercent{}, or about \idvPeakErrorKelvin. \par
%
\begin{table}[tb]
    \centering
    \caption{Parameters of \texttt{IDV(29)}.}
    \label{tab:idv29}
    \begin{tabular}{lll}
        \toprule
        Symbol & Description & Value \\
        \midrule
        $\kappa_m$ & Drift coefficient & \idvDriftCoefficient \\
        $T_{\mathrm{seg}}$ & Nominal segment duration & \idvSegmentDuration \\
        $\gamma_\tau$ & Relative shift of the segment boundaries & \idvShiftFraction \\
        $c_i^{+}$ & Start level, drawn uniformly from & $[\idvUpperBandLow, \idvUpperBandHigh]$ \\
        $c_i^{-}$ & Target level, drawn uniformly from & $[\idvLowerBandLow, \idvLowerBandHigh]$ \\
        \bottomrule
    \end{tabular}
\end{table}
% 
Since the Simulink model integrates the plant with the solver \texttt{ode45}, whose step size varies, the model function is called at points in time that do not lie on a fixed grid. The drift signal is therefore formulated as a closed-form expression in the simulation time $t$ and evaluated anew at every call. The displacement of each boundary and the two levels of each segment are derived deterministically from a seed and the segment index, so that the signal carries no state from one call to the next. The signal is thus independent of the step size the solver chooses, and repeated runs with the same seed carry an identical drift. Only the number of segments follows from the simulated duration. \par
%
%
The intervention in the source code of the simulator is confined to a few lines of \texttt{temexd\_mod.c}, shown in listing~\ref{lst:idv29}. The drift signal itself is implemented in the separate file \texttt{drift\_function.c}, which the S-function includes as source and which can also be compiled into a shared library, so that the identical signal can be evaluated from Python. In the listing, \texttt{tcc} denotes the stripper temperature computed by the process model, which is first assigned to the zero-indexed measurement vector \texttt{xmeas} at the position of \texttt{XMEAS(18)}. The entry \texttt{idv[29]} holds the value the disturbance is set to through the interface, which is the drift coefficient $\kappa_m$. On the first call after activation, the simulation time \texttt{time} is stored in \texttt{drift\_tonset}, which is reset to a negative value once the disturbance is switched off, so that a renewed activation restarts the pattern. The function \texttt{drift\_signal} returns $c(t)$ for the time elapsed since activation and is parameterized by the constants \texttt{DRIFT\_PERIOD}, \texttt{DRIFT\_GAMMA} and \texttt{DRIFT\_SEED}, which correspond to $T_{\mathrm{seg}}$, $\gamma_\tau$ and the seed of the random levels. The resulting factor \texttt{driftf} scales the entry at \texttt{DRIFT\_SENSOR\_INDEX}, which refers to \texttt{XMEAS(18)}, only after the process value has been assigned. It therefore acts solely on the recorded measurement, while the simulated process behaves exactly as in the undisturbed case. Both files, together with the adjustments to \texttt{pyTEP} required to run the extended simulation, are publicly available in a fork of the package~\cite{marijan.2026pytep}. \par
%
\begin{lstlisting}[language=C, float=tb,
    caption={Application of the sensor drift in \texttt{temexd\_mod.c}, abridged and simplified.},
    label={lst:idv29}]
xmeas[17] = tcc;                      /* process value of XMEAS(18) */
if (idv[29] > 0.) {
    if (drift_tonset < 0.) drift_tonset = time;
    driftf = 1. + idv[29] * drift_signal(time - drift_tonset,
                              DRIFT_PERIOD, DRIFT_GAMMA, DRIFT_SEED);
} else {
    drift_tonset = -1.;
    driftf = 1.;
}
xmeas[DRIFT_SENSOR_INDEX] *= driftf;
\end{lstlisting}